\section{Why these clauses are well-defined}
The original draft raised two technical problems: division by zero in the residual definitions and
the treatment of words not present in the corpus. The present formalisation resolves both.

First, division by zero is avoided by restricting the residual minima to witnesses with positive
denominator, i.e. to contexts $y$ such that $v(y,A)>0$. If no such witness exists, the set in the
minimum collapses to $\{1\}$ and the value becomes $1$. This is a natural default: in the absence
of any positively supported counterevidence, the residual connective is not forced below $1$.
The inner truncation $\min(1,\_)$ ensures that the value stays in the interval $[0,1]$.

Second, out-of-vocabulary words are excluded from the domain entirely. The semantics is not a
universal language model over all possible strings; it is a finite, corpus-relative semantics for
observed fragments. This is consistent with the stated goal of the thesis, namely to parse strings
of words drawn from the corpus vocabulary.

\begin{remark}[Subtleties not captured by the present formalisation]
The following issues are intentionally left outside the formal model.
\begin{enumerate}[label=(\alph*)]
    \item \textbf{Learning from raw corpora.} The formalisation assumes lexical probabilities are
    already given, whereas a practical system would need to infer them.
    \item \textbf{Out-of-vocabulary generalisation.} Smoothing, backoff, or subword modelling would be
    needed for real applications.
    \item \textbf{Global normalisation.} The values $v(x,A)$ behave as probabilities locally, but the model is
    not presented here as a full generative probability distribution over derivation trees.
    \item \textbf{Unseen but grammatical strings.} By working over $\Frag(\mathcal C)$, the semantics is tied to
    attested fragments rather than all strings over $\Sigma$.
\end{enumerate}
These restrictions should be read as part of the thesis scope, not as permanent limitations on
future versions of the model.
\end{remark}

